%! $n$th Roots \& Rational Exponents — Unit 6, Lesson 6.1 — Student Packet Cover Sheet
\documentclass[10pt]{article}
\usepackage{algebra2-article}
\usepackage{algebra2-boxes}
\usepackage{ltablex}
\keepXColumns

\begin{document}

% Full-bleed forest banner
\begin{tikzpicture}[remember picture, overlay]
  \fill[forest] (current page.north west)
    rectangle ([yshift=-0.9in]current page.north east);
\end{tikzpicture}
\vspace{-0.8in}

\begin{center}
  {\color{white}\textbf{\LARGE Algebra 2: Shepherd}}\\[4pt]
  {\color{forestmid}\large\textbf{Unit 6 \quad Radical Functions}}\\[6pt]
  {\color{white}\textbf{\large Lesson 6.1 \quad $n$th Roots \& Rational Exponents}}
\end{center}

\vspace{0.22in}
\namedateperiod
\vspace{0.18in}

\begin{learningtargetbox}
\small
\begin{enumerate}[leftmargin=1.5em, itemsep=5pt]
  \item \ldots{} find \textbf{$n$th roots}, name the \textbf{principal} root, and explain why an \textbf{even} root of a negative number is not a real number while an \textbf{odd} one is.
  \item \ldots{} \textbf{convert} between radical form and \textbf{rational exponent} form using $a^{m/n}=\sqrt[n]{a^m}=\left(\sqrt[n]{a}\right)^m$, and \textbf{evaluate} such expressions without a calculator.
  \item \ldots{} apply the \textbf{exponent properties} to rational exponents --- including combining radicals with \textbf{different indices} --- and use the denominator to decide whether a domain is restricted.
\end{enumerate}
\end{learningtargetbox}

\vspace{0.14in}

\begin{tocbox}
\small
\renewcommand{\arraystretch}{1.5}
\begin{tabularx}{\linewidth}{@{} c l X r @{}}
\rowcolor{forestbg}
\textbf{\#} & \textbf{Component} & \textbf{Description} & \textbf{Score} \\
\midrule
1 & Warm-Up        & Spiral review: perfect powers, the exponent properties, and fraction arithmetic & \blank{1.2cm} \\
2 & Guided Notes   & $n$th roots and principal roots, $\sqrt[n]{a}=a^{1/n}$, the $a^{m/n}$ rule, and the properties applied to radicals & \blank{1.2cm} \\
3 & Group Activity & \emph{One Expression, Two Costumes} --- translating, evaluating, and a Kepler's-law model, by tier & \blank{1.2cm} \\
4 & Exit Ticket    & Quick check: convert, evaluate, explain even vs.\ odd, plus an SOL-style item & \blank{1.2cm} \\
5 & Homework       & Mixed conversion and evaluation, the absolute-value rule, domains, and a metabolic-rate model & \blank{1.2cm} \\
\midrule
  & \multicolumn{2}{r}{\textbf{Total}} & \blank{1.2cm} \\
\end{tabularx}
\end{tocbox}

\vspace{0.10in}

\begin{remindbox}
\small A radical is an \textbf{exponent in disguise}. Nobody decided this --- the exponent rules
force it: if ``power to a power'' is to keep working, then $\left(a^{1/n}\right)^{n}=a^{1}=a$, so
$a^{1/n}$ must be the number whose $n$th power is $a$. That is the definition of
$\sqrt[n]{a}$. Hence $\sqrt[n]{a}=a^{1/n}$ and, more generally,
$a^{m/n}=\sqrt[n]{a^{m}}=\left(\sqrt[n]{a}\right)^{m}$: the \textbf{denominator} is the
\textbf{index} (the root you take) and the \textbf{numerator} is the \textbf{power}. Flipping those
two is the most common mistake in this lesson --- $\sqrt[4]{x^{3}}$ is $x^{3/4}$, never $x^{4/3}$.
When you evaluate, take the \textbf{root first}: $16^{3/4}=\left(\sqrt[4]{16}\right)^{3}=2^{3}=8$ is
mental arithmetic, while $\sqrt[4]{16^{3}}=\sqrt[4]{4096}$ is not. A \textbf{negative} rational
exponent means \emph{reciprocal}, never a negative answer: $16^{-1/2}=\frac{1}{4}$. And the payoff:
because these are just exponents, every exponent property you already know now applies to radicals
--- including ones no radical rule can touch, like
$\sqrt{x}\cdot\sqrt[3]{x}=x^{1/2}\cdot x^{1/3}=x^{5/6}=\sqrt[6]{x^{5}}$.
\end{remindbox}

\end{document}
